\subsection{The binary affine family}\label{the-binary-affine-family}

Let \(S_d\) consist of points and affine planes of AG(d,2), d≥3. Every
three distinct points are affinely independent over F₂ and lie in a
unique affine plane, so this is \(S(3,4,2^d)\). Deriving at the origin
identifies the remaining points with nonzero vectors. A derived block is
\{u,v,u+v\}, precisely a projective line in PG(d−1,2). Thus the derived
design is PG(d−1,2).

The projective base theorem above and Theorem S.5 of the common paper
show: \[
\Gamma(S_d)\text{ is residually connected}\iff d=3.
\tag{C4}
\] This proves L092 Lemma4 without repeating an obstruction at every
rank. The positive base PG(2,2) is included explicitly; the source's
Theorem3 excludes q=2, so one must not cite that restricted statement
alone for this application.

The full affine group is chamber-transitive for every d≥3. To apply
Theorem S.3, note first that AGL(d,2) is transitive on ordered triples
of distinct points: translate the first point to zero, then send the two
independent difference vectors to prescribed independent vectors.
Normalize a triple to (0,e₁,e₂). The unique block on it is ⟨e₁,e₂⟩.
Alternative planes through the line ⟨e₁⟩ correspond to nonzero vectors
in \[
\mathbb F_2^d/\langle e_1\rangle
\] other than the vector represented by e₂. The stabilizer fixing e₁,e₂
induces the full stabilizer of that nonzero vector in GL(d−1,2). This
group is transitive on the remaining nonzero vectors: any ordered pair
of independent vectors extends to a basis. Theorem S.3 proves chamber
transitivity.

At d=3 the type sizes are 8 ordinary points, 84 marked objects, 14
blocks and 8 copied points; the panel sizes are 2,3,2,2 in that type
order. There are (8)₃·2=672 chambers. The group AGL(3,2), of order 1344,
has chamber stabilizer of order 2. These formulas give the numerical
positive case in Theorem 11 of {[}L{]}. The full design group is exactly
AGL(d,2). Indeed, after translating the image of zero back to zero, a
design automorphism preserves the unique fourth point x+y of the affine
plane through 0,x,y whenever x,y are distinct nonzero vectors. It
therefore preserves addition in this case; equal vectors or a zero
vector give the same identity immediately in characteristic two. Hence
it is linear. This proves the full-group identification directly, while
chamber transitivity already follows from the explicit subgroup
exhibited above.

\subsection{The Mathieu applications}\label{the-mathieu-applications}

For S(4,5,11) and S(5,6,12), the chamber formula of Theorem S.3 gives
respectively 23,760 and 285,120 chambers. The literal finite computation
\texttt{phase186\_witt\_derived\_check.g} constructs their forced native block
systems, verifies the Steiner property and computes full automorphism
groups of orders 7,920 and95,040. Neither group can therefore be
transitive: the chamber count is three times the full group order. No
search over putative chamber orbits is necessary.

For S(3,6,22), S(4,7,23) and S(5,8,24), the common value is r=5 and the
final fiber has size 4. Their chamber counts are
respectively 36,960,850,080 and20,401,920. The complete program
\texttt{phase182\_mathieu\_fibers.g} constructs the actual block
designs, verifies the Steiner property over every t-subset, and checks
the stabilizer orbits supplying the source's Property(*). That
computation is retained; it is not duplicated here.

Their (t−2)-fold derived design is PG(2,4). This is now established by
explicit coordinate-incidence maps, described below and checked by GAP,
independent coefficient arithmetic, and the finite Lean certificate
\texttt{WittDerivedPlanes.lean}. The bridge from the original native
Witt blocks to each displayed derived incidence table remains checked in
GAP and Python. The residue recursion, Theorem S.6 and the projective
base theorem then prove geometry, firmness, residual connectivity and
(IP)₂, while the retained fiber computation proves chamber transitivity.

For the rank-six M24 construction, the ordinary-point residue is
Γ(S(4,7,23)), with type sizes 23,253,8,855,253,23. The residue of an
octad block has type sizes 8,28,56,280,16: its ordinary subsets have
sizes 1,2,3; a marked object chooses a four-subset of the octad and one
of four alternative octads; and its copied points lie outside the octad.
These two type-size multisets prove that the displayed rank-five
residues are not isomorphic. The Figure 3 subgroup labels are proved
below by explicit finite kernels, complements and transported actions.
The object-level identifications with the named earlier Pasini
geometries remain source-specific obligations.
